\documentclass[12pt]{article}
\usepackage[utf8]{inputenc}
\usepackage[T1]{fontenc}
\usepackage{hyperref}
\title{Sherdukpen Language Resources for K-12 Curriculum Development}
\author{Compiled by Claude Sonnet 5.5 (claude-sonnet-5-5) for Andr\'as Kornai}
\date{}
\begin{document}
\maketitle

\section{Introduction}
Sherdukpen (autonym Mey; ISO 639-3 sdp) is a Kho-Bwa Sino-Tibetan language of West Kameng district, Arunachal Pradesh, India. It has two varieties, Rupa and Shergaon. Published estimates of speaker numbers differ. The language is not well documented, and no school materials were found.

\section{Primary Educational Resources}
No primers or textbooks were found. A news report describes a project to develop a Roman-script alphabet, dictionary and grammar. A community leader is quoted as saying Sherdukpen would then be introduced as a third language in schools in Sherdukpen areas.
\begin{itemize}
\item Arunachal Times, ``Written script for Sherdukpen language underway''. \url{https://arunachaltimes.in/?p=252775}
\end{itemize}

\section{Linguistic and Reference Materials}
\begin{itemize}
\item Wikipedia article on the Sherdukpen language. \url{https://en.wikipedia.org/wiki/Sherdukpen_language}
\item eScholarship item (Himalayan Linguistics) containing a phonological study of the Rupa variety. \url{https://escholarship.org/uc/item/3bc1m88z}
\item Himalayan Linguistics, volume 23 issue 2 listing. \url{https://escholarship.org/uc/himalayanlinguistics/23/2}
\item Brahmaputra Studies fieldwork page on Sherdukpen (Dollfus and Jacquesson, CNRS). \url{https://brahmaputra.ups2259.vjf.cnrs.fr/uk/fieldworks/Sherdukpen/Sherdukpen1.htm}
\item ASJP word list for Sherdukpen. \url{https://asjp.clld.org/languages/SHERDUKPEN}
\end{itemize}
A 1988 handbook on Sherdukpen by Rinchin Dondrup (Directorate of Research, Government of Arunachal Pradesh) is cited in secondary sources. No online copy or URL was found.

\section{Religious and Community Resources}
\begin{itemize}
\item Wikipedia article on the Sherdukpen people. \url{https://en.wikipedia.org/wiki/Sherdukpen_people}
\item Asia Harvest people-group profile. It classifies the language differently from the Kho-Bwa placement used elsewhere. \url{https://www.people-groups.asiaharvest.org/India/Sherdukpen.pdf}
\item Joshua Project language page for sdp. \url{https://joshuaproject.net/index.php/languages/sdp}
\end{itemize}

\section{Digital and Audio Resources}
The CIIL Bhasha Sanchika special collection came up in search results. Its Sherdukpen content was not opened or confirmed.
\begin{itemize}
\item \url{https://sanchika.ciil.org/communities/5cd82e96-f8ef-4a80-927e-b5325b6c9662}
\end{itemize}

\section{Curriculum Development Context}
The language is documented only thinly and has no standardised script yet. Any curriculum would first need an agreed orthography, which the Shergaon script project may provide. It would also need a choice between the Rupa and Shergaon varieties, and primary materials produced with community speakers.
\end{document}